% !TeX root = ../drafts/08-end-to-end-protocol.tex
\section{End-to-end protocol}\label{sec:e2e}

\subsection{Bytecode encoding}\label{sec:e2e-bc}

The bytecode (of $\nprog=2^{\kbc}$ instructions) is encoded as one $\K$-valued multilinear $P$ in $\kbc+4$ variables:
\begin{enumerate}[itemsep=1pt,topsep=3pt]
\item lay instruction $z$ out as $16$ slots of $\K$: the opcode in slot $3$, the operands and immediate lanes in slots $4,\dots,11$, and $0$ everywhere else.
\item read the resulting $2^{\kbc}\times16$ table as $P$, the low $\kbc$ variables indexing the instruction and the high four its slot, bits low first.
\end{enumerate}
\begin{center}
\begin{tabular}{llllllllll}
\hline
Instruction & $3$ & $4$ & $5$ & $6$ & $7$ & $8$ & $9$ & $10$ & $11$\\
\hline
\texttt{XOR}           & $\opc{XOR}$ & $o_A$ & $o_B$ & $o_C$ & $0$ & $0$ & $0$ & $0$ & $0$\\
\texttt{MUL\_NATIVE}   & $\opc{MUL}$ & $o_A$ & $o_B$ & $o_C$ & $0$ & $0$ & $0$ & $0$ & $0$\\
\texttt{SET\_CONSTANT} & $\opc{SET}$ & $o$   & $k_0$ & $k_1$ & $k_2$ & $0$ & $0$ & $0$ & $0$\\
\texttt{DEREF}         & $\opc{DRF}$ & $o_1$ & $o_2$ & $o_3$ & $f_{\pc}$ & $f_{\fp}$ & $0$ & $0$ & $0$\\
\texttt{JUMP}          & $\opc{JMP}$ & $o_c$ & $o_d$ & $o_f$ & $0$ & $0$ & $0$ & $0$ & $0$\\
\texttt{BLAKE2S}        & $\opc{B2S}$  & $o_{m_0}$ & $o_{m_1}$ & $o_{m_2}$ & $o_{m_3}$ & $o_{\mathit{cv}}$ & $o_{\mathit{out}}$ & $\md_0$ & $\md_1$\\
\hline
\end{tabular}
\end{center}
So $P(z,1,1,0,0)$ is instruction $z$'s opcode, $P(z,0,0,1,0)$ its first operand, and $P(z,0,0,0,0)=P(z,0,0,1,1)=0$.

\subsection{Public input}\label{sec:e2e-pi}

The public input is the first two memory cells $\mem[\gen^{0}],\mem[\gen^{1}]$, two $192$-bit words (with top limb $0$) both parties know. The verifier samples $r_m\in\E$, the prover sends $c_0,c_1$ claiming $c_\ell=\mle{\mem_\ell}(r_m,0,\dots,0)$, and the verifier checks
\[
  c_\ell\;\qeq\;(1+r_m)\,\mem[\gen^{0}]_\ell+r_m\,\mem[\gen^{1}]_\ell,\qquad \ell\in\{0,1\}.
\]
Both sides are degree one in $r_m$, so a limb disagreeing with the public words survives with probability at most $1/|\E|$ (Lemma~\ref{lem:sz}). The claims $c_0$, $c_1$ are discharged to the PCS.

\subsection{Filling the tables}\label{sec:e2e-pad}

A table is proven over a power-of-two number of rows. The prover reaches one by running extra instructions, in execution cycles disjoint from the program's own run. Proposition~\ref{prop:walk} about bus balance splits the rows into two groups:
\begin{itemize}
\item one walk from $(\pc_0,\fp_0)$ to $(\pc_{\mathrm{final}},\fp_{\mathrm{final}})$
\item closed walks
\end{itemize}
The prover can thus leverage these \emph{closed walks} to fill every table to its next power of 2. (important for completeness, but not for soundness)

\subsection{The unrolled protocol}\label{sec:e2e-unrolled}

Each reduction below adds evaluation claims on columns to a shared \emph{claim pool}; the final \textsc{Opening} phase proves all of them with a single PCS opening.

\paragraph{Setup and statement binding.} Fixed and shared in advance through \S\ref{sec:e2e-pi}): the fields and constants; the instance caps (\S\ref{sec:lookup}); the bytecode and its length; the initial state $(\pc_0,\fp_0)$ and final state $(\pc_{\mathrm{final}},\fp_{\mathrm{final}})$; and the public input $\mem[\gen^{0}],\mem[\gen^{1}]$. Before any challenge, the transcript is seeded by the public input and one environment digest that binds both the exact bytecode and flock's BLAKE2s R1CS. The prover then announces the memory log-size $\kmem$ and the six table log-heights $\tau_j$; the verifier checks them against the caps and derives every table, stack, and leaf shape.

\paragraph{Commitment.}
\begin{enumerate}
\item \textbf{Prover.} Commits every column: each table's state, operands, read values and counts; the three memory limbs; the two finalize counts; and Flock's packed witness $\qflock$ (\S\ref{sec:tab-blake2s}). A program that never executes \texttt{BLAKE2S} still runs Flock's least instance count of dummy compressions, so the proof shape stays uniform. The prover stacks the columns into $q$ and sends the commitment (\S\ref{sec:stacking}, Annex~\ref{annex:pcs}).
\item \textbf{Verifier.} Checks the announced sizes against the caps and derives the table, stack and leaf shapes (\S\ref{sec:leafstack}).
\end{enumerate}

\paragraph{Bus.}
\begin{enumerate}
\item \textbf{Verifier.} Sends the fingerprint challenge $\alpha\in\E^{4}$ and the multiset challenge $\beta\in\E$ (\S\ref{sec:gp}).
\item \textbf{Prover.} Forms the push, pull, and \emph{count} leaf vectors (\S\ref{sec:leafstack}; the count leaves are the count columns themselves, \S\ref{sec:lookup}), stacked into blocks and padded with $1$s. Sends the count root $R_c$ and one bus root $R$, the value both the push and the pull product must take (\S\ref{sec:gp}).
\item \textbf{Prover \& Verifier.} Run the one batched GKR (\S\ref{sec:gkr}) over all three trees. For the bus, the verifier checks each side's product against $R$; one $R$ for both sides is the balance check (\S\ref{sec:gp}). For the counts it checks $R_c\neq0$ (no read self-cancels, \S\ref{sec:lookup}). The pass yields three leaf claims $\widetilde V^{s}_0(\zeta)$ at one shared $\zeta$.
\item \textbf{Prover \& Verifier.} Decompose each leaf claim (\S\ref{sec:leafstack}). The verifier forms every public part itself: the block selectors, the constant entries, the index column (\S\ref{sec:idxcol}), and the program's whole share of the two bytecode blocks, one evaluation at $(\zeta,\alpha)$ (\S\ref{sec:e2e-bc}). Three blocks per side belong to no table: the state boundary, and each array's seed and finalization. For those the prover sends one value per committed column, which enter the claim pool. Every other block belongs to a table, and for those it sends nothing: what they owe is the leaf claim less the part the verifier just formed, and the table sumcheck settles it (\S\ref{sec:air}).
\end{enumerate}

\paragraph{Row sumcheck.} One run for the six tables' constraints and their bus forms (\S\ref{sec:air}).
\begin{enumerate}
\item \textbf{Verifier.} Sends the batch challenge $\xi\in\E$. Each table takes a disjoint range of its powers, one per constraint, and the $\nside$ bus forms take one power each, shared by every table. No point is drawn: the batch runs at the bus point, table $T_j$ tested at $\zeta_{<\tau_j}$ (\S\ref{sec:air}).
\item \textbf{Prover \& Verifier.} Run the table sumcheck, $\taumax=\max_j\tau_j$ rounds, a table of $2^{\tau}$ rows joining at round $\taumax-\tau$. The target is the one the verifier derived from the three leaf claims. A round message is a degree-three univariate, and the round check pins one of its four coefficients (\S\ref{sec:prelim}), so each round costs three $\E$-elements. The end yields one claim per column of every table, at that table's own point $\fc_{<\tau_j}$; these enter the pool, and the bus forms are read off them. \texttt{BLAKE2S}'s value columns are not committed here: their claims route to $\qflock$ (\S\ref{sec:tab-blake2s}).
\end{enumerate}

\paragraph{Public input.}
\begin{enumerate}
\item \textbf{Verifier.} Sends $r_m\in\E$ and forms the line $\widetilde{\mathrm{PI}}_\ell(r_m)$ from the public words.
\item \textbf{Prover \& Verifier.} The prover sends $c_0,c_1$; the verifier checks them against the line and pools them, with $0$ for the top limb, as claims on $\mem_0,\mem_1,\mem_2$ at $(r_m,0,\dots,0)$ (\S\ref{sec:e2e-pi}).
\end{enumerate}

\paragraph{BLAKE2s validity.}
\begin{enumerate}
\item \textbf{Prover \& Verifier.} Reduce the executed BLAKE2s constraints to one evaluation claim on $\qflock$ (\S\ref{sec:tab-blake2s}, Annex~\ref{annex:flock}).
\item \textbf{Prover \& Verifier.} The ring switching of \S\ref{sec:ringswitch} reduces that bit-witness claim to a weighted-sum claim on the $\qflock$ region of the stack, with an MLE-friendly $\E$-valued weight (Definition~\ref{def:ipcs}); it joins the pool.
\end{enumerate}

\paragraph{Opening.}
\begin{enumerate}
\item \textbf{Prover \& Verifier.} Each pooled claim is one evaluation of a column at a point, its value already supplied; via the stacking selectors (\S\ref{sec:stacking}) it becomes a weighted sum $\sum_w W_j(w)\,\widetilde q(w)=c_j$ over the stacked witness. (The ring-switched claim arrives already in this form, its weight supported on the $\qflock$ region.)
\item \textbf{Verifier.} Sends one batching challenge $\lambda\in\E$, drawn after every claim value is bound. Claim $j$ takes the weight $\lambda^{j-1}$, the ring-switched claim first and the pooled column claims after: the $J$ weights are distinct powers of the one challenge (Theorem~\ref{thm:rbr}).
\item \textbf{Prover \& Verifier.} Fold the $J$ claims into $W_\lambda=\sum_j\lambda^{j-1} W_j$ and $C_\lambda=\sum_j\lambda^{j-1} c_j$, and run the inner-product PCS opening on $\sum_w W_\lambda(w)\,\widetilde q(w)=C_\lambda$ (Annex~\ref{annex:pcs}), the verifier evaluating $W_\lambda$ itself. This one opening discharges every pooled claim; there is no separate reduction sumcheck.
\item \textbf{Verifier.} Accepts if every check above passed: the caps, the one bus root for both sides, the nonzero count root, every sumcheck's rounds and final value, the public-input line, flock's reduction, and the PCS opening.
\end{enumerate}
